% GENERATED FILE. Edit canonical lessons, then run npm run generate:books.
% canonical-source-hash: fnv1a64:17d63ca6d240bd1b
% canonical-lessons: FA-C523-teatr, FA-C523-sahne, FA-C523-namayeshname, FA-C523-tamashachi, FA-C523-naqsh

\chapter{At the Theatre}
\label{ch:fa-at-the-theatre}
\hlchaptermodality{523}

\begin{chapteropening}
\textbf{By the end of this chapter:} \emph{I can talk about going to the theatre.}

Complete the last lesson of chapter 523: I can talk about going to the theatre.
\end{chapteropening}

\section[\texorpdfstring{\fa{تئاتر} (te'âtr)}{تئاتر (te'âtr)}]{\fa{تئاتر} (te'âtr) — a theatre; a play}

\label{lesson:FA-C523-teatr}



\begin{quote}
Before the new one: say the Persian for to be inspired, then the Persian for imagination.
\end{quote}

\subsection*{You'll want to know: \fa{تئاتر}}
\textbf{\fa{تئاتر}} (\emph{te'âtr}) — \textquotedblleft{}a theatre; a play\textquotedblright{}.

\begin{quote}
\fa{ما} \fa{دیشب} \fa{به} \fa{تئاتر} \fa{رفتیم}.
\end{quote}

\begin{quote}
\emph{mâ dishab be te'âtr raftim.}
\end{quote}

\begin{quote}
We went to the theatre last night.
\end{quote}

Both the place and the play.

\begin{cousinweb}
\emph{te'âtr} was borrowed from French \emph{théâtre}, which EtymDB traces through Latin \emph{theatrum} to Greek \emph{theatron}, 'a place for viewing' (EtymDB); Lewis \& Short gloss \emph{theatrum} 'a theatre' (Lewis \& Short).
\end{cousinweb}

\begin{grammarlens}[title={What you've built}]
One more word for telling, explaining and describing at B1.
\end{grammarlens}

\subsection*{Your turn}
\begin{itemize}
\raggedright
  \item \emph{Say it:} \emph{te'âtr}
  \item \emph{Say it:} \emph{te'âtr}, once more
  \item \emph{Say it:} a sentence of your own with \emph{te'âtr}
  \item \emph{Recall:} say the Persian for to be inspired, then the Persian for imagination, then say \emph{te'âtr} again
\end{itemize}

\begin{tcolorbox}[breakable,colback=teal!4,colframe=teal!35!black,title={Before you move on}]
What does \fa{تئاتر} mean? (\textquotedblleft{}A theatre; a play\textquotedblright{}.) Say it once more.
\end{tcolorbox}
\section[\texorpdfstring{\fa{صحنه} (sahne)}{صحنه (sahne)}]{\fa{صحنه} (sahne) — a stage; a scene}

\label{lesson:FA-C523-sahne}



\begin{quote}
Before the new one: say the Persian for imagination, then the Persian for a theatre.
\end{quote}

\subsection*{You'll want to know: \fa{صحنه}}
\textbf{\fa{صحنه}} (\emph{sahne}) — \textquotedblleft{}a stage; a scene\textquotedblright{}.

\begin{quote}
\fa{بازیگرها} \fa{روی} \fa{صحنه} \fa{آمدند}.
\end{quote}

\begin{quote}
\emph{bâzigarhâ ru-ye sahne âmadand.}
\end{quote}

\begin{quote}
The actors came on stage.
\end{quote}

Also a scene in a film or a play. Plural \emph{sahnehâ}.

\begin{cousinweb}
\emph{sahne} is a loan, usually traced to Arabic \emph{ṣaḥn}, 'courtyard, open space'. EtymDB lists it without a source.
\end{cousinweb}

\begin{grammarlens}[title={What you've built}]
One more word for telling, explaining and describing at B1.
\end{grammarlens}

\subsection*{Your turn}
\begin{itemize}
\raggedright
  \item \emph{Say it:} \emph{sahne}
  \item \emph{Say it:} \emph{sahne}, once more
  \item \emph{Say it:} a sentence of your own with \emph{sahne}
  \item \emph{Recall:} say the Persian for imagination, then the Persian for a theatre, then say \emph{sahne} again
\end{itemize}

\begin{tcolorbox}[breakable,colback=teal!4,colframe=teal!35!black,title={Before you move on}]
What does \fa{صحنه} mean? (\textquotedblleft{}A stage; a scene\textquotedblright{}.) Say it once more.
\end{tcolorbox}
\section[\texorpdfstring{\fa{نمایشنامه} (namâyeshnâme)}{نمایشنامه (namâyeshnâme)}]{\fa{نمایشنامه} (namâyeshnâme) — a play (script)}

\label{lesson:FA-C523-namayeshname}



\begin{quote}
Before the new one: say the Persian for a theatre, then the Persian for a stage.
\end{quote}

\subsection*{You'll want to know: \fa{نمایشنامه}}
\textbf{\fa{نمایشنامه}} (\emph{namâyeshnâme}) — \textquotedblleft{}a play (script)\textquotedblright{}.

\begin{quote}
\fa{او} \fa{یک} \fa{نمایشنامه} \fa{کوتاه} \fa{برای} \fa{بچه} \fa{ها} \fa{نوشت}.
\end{quote}

\begin{quote}
\emph{u yek namâyeshnâme-ye kutâh barâ-ye bachchehâ nevesht.}
\end{quote}

\begin{quote}
She wrote a short play for children.
\end{quote}

\emph{namâyesh} 'show' + \emph{nâme} 'writing, letter'. Add \emph{-nevis} 'writer' for a playwright.

\begin{cousinweb}
\emph{namâyeshnâme} is built on native Persian parts: \emph{namâyesh}, 'show', from the verb \emph{namudan}, 'to show', and \emph{nâme}, 'letter, writing'. EtymDB has no entry for the compound.
\end{cousinweb}

\begin{grammarlens}[title={What you've built}]
One more word for telling, explaining and describing at B1.
\end{grammarlens}

\subsection*{Your turn}
\begin{itemize}
\raggedright
  \item \emph{Say it:} \emph{namâyeshnâme}
  \item \emph{Say it:} \emph{namâyeshnâme}, once more
  \item \emph{Say it:} a sentence of your own with \emph{namâyeshnâme}
  \item \emph{Recall:} say the Persian for a theatre, then the Persian for a stage, then say \emph{namâyeshnâme} again
\end{itemize}

\begin{tcolorbox}[breakable,colback=teal!4,colframe=teal!35!black,title={Before you move on}]
What does \fa{نمایشنامه} mean? (\textquotedblleft{}A play (script)\textquotedblright{}.) Say it once more.
\end{tcolorbox}
\section[\texorpdfstring{\fa{تماشاچی} (tamâshâchi)}{تماشاچی (tamâshâchi)}]{\fa{تماشاچی} (tamâshâchi) — a spectator, a member of the audience}

\label{lesson:FA-C523-tamashachi}



\begin{quote}
Before the new one: say the Persian for a stage, then the Persian for a play.
\end{quote}

\subsection*{You'll want to know: \fa{تماشاچی}}
\textbf{\fa{تماشاچی}} (\emph{tamâshâchi}) — \textquotedblleft{}a spectator, a member of the audience\textquotedblright{}.

\begin{quote}
\fa{تماشاچی} \fa{ها} \fa{برای} \fa{بازیگرها} \fa{دست} \fa{زدند}.
\end{quote}

\begin{quote}
\emph{tamâshâchihâ barâ-ye bâzigarhâ dast zadand.}
\end{quote}

\begin{quote}
The audience applauded the actors.
\end{quote}

\emph{tamâshâ} 'watching' + the Turkic ending \emph{-chi} 'one who'.

\begin{cousinweb}
\emph{tamâshâchi} is built in Persian on \emph{tamâshâ}, 'watching', a loan from Arabic, and the ending \emph{-chi}, 'one who', which was borrowed from Turkic. EtymDB has no entry for it.
\end{cousinweb}

\begin{grammarlens}[title={What you've built}]
One more word for telling, explaining and describing at B1.
\end{grammarlens}

\subsection*{Your turn}
\begin{itemize}
\raggedright
  \item \emph{Say it:} \emph{tamâshâchi}
  \item \emph{Say it:} \emph{tamâshâchi}, once more
  \item \emph{Say it:} a sentence of your own with \emph{tamâshâchi}
  \item \emph{Recall:} say the Persian for a stage, then the Persian for a play, then say \emph{tamâshâchi} again
\end{itemize}

\begin{tcolorbox}[breakable,colback=teal!4,colframe=teal!35!black,title={Before you move on}]
What does \fa{تماشاچی} mean? (\textquotedblleft{}A spectator, a member of the audience\textquotedblright{}.) Say it once more.
\end{tcolorbox}
\section[\texorpdfstring{\fa{نقش} (naqsh)}{نقش (naqsh)}]{\fa{نقش} (naqsh) — a role, a part; a pattern}

\label{lesson:FA-C523-naqsh}



\begin{quote}
Before the new one: say the Persian for a play, then the Persian for a spectator.
\end{quote}

\subsection*{You'll want to know: \fa{نقش}}
\textbf{\fa{نقش}} (\emph{naqsh}) — \textquotedblleft{}a role, a part; a pattern\textquotedblright{}.

\begin{quote}
\fa{او} \fa{در} \fa{این} \fa{فیلم} \fa{نقش} \fa{یک} \fa{پدر} \fa{را} \fa{دارد}.
\end{quote}

\begin{quote}
\emph{u dar in film naqsh-e yek pedar râ dârad.}
\end{quote}

\begin{quote}
He plays a father in this film.
\end{quote}

Also 'design, pattern' on a carpet or a tile.

\begin{cousinweb}
\emph{naqsh} is a loan, usually traced to Arabic \emph{naqsh}, 'a painting, an engraving'; Persian widened it to a pattern and then to a part in a play. EtymDB has no entry for it.
\end{cousinweb}

\begin{grammarlens}[title={What you've built}]
One more word for telling, explaining and describing at B1.
\end{grammarlens}

\subsection*{Your turn}
\begin{itemize}
\raggedright
  \item \emph{Say it:} \emph{naqsh}
  \item \emph{Say it:} \emph{naqsh}, once more
  \item \emph{Say it:} a sentence of your own with \emph{naqsh}
  \item \emph{Recall:} say the Persian for a play, then the Persian for a spectator, then say \emph{naqsh} again
\end{itemize}

\begin{tcolorbox}[breakable,colback=teal!4,colframe=teal!35!black,title={Before you move on}]
What does \fa{نقش} mean? (\textquotedblleft{}A role, a part; a pattern\textquotedblright{}.) Say it once more.
\end{tcolorbox}
